%% kapa: exam=1 round=24 year=2013 period=2 form=A subject=accounting subject_ko=회계학 seq=26 number=26
%% source: Q-Net 감정평가사 기출문제 (cst00309) artlSeq=5073802 file=extracted/1008/5073802/2013년 제24회 감정평가사 1차시험 기출문제/2교시(A형).hwp
%% conversion: kordoc 4.24.51 -> tools/gen1.py (auto); review=proofread-sonnet; flags=none
\begin{question}{26}
주식기준보상에 관한 설명으로 옳지 않은 것은?
\begin{choices}
  \item 현금결제형 주식기준보상거래의 경우, 제공받는 재화나 용역과 그 대가로 부담하는 부채를 부채의 공정가치로 측정한다.
  \item 현금결제형 주식기준보상거래의 경우, 부채가 결제될 때까지 매 보고기간말과 결제일에 부채의 공정가치를 재측정하고, 공정가치의 변동액은 기타포괄손익으로 인식한다.
  \item 주식결제형 주식기준보상거래의 경우, 제공받는 용역의 공정가치를 신뢰성 있게 추정할 수 없다면, 제공받는 용역과 그에 상응하는 자본의 증가는 부여된 지분상품의 공정가치에 기초하여 간접 측정한다.
  \item 주식결제형 주식기준보상거래의 경우, 부여한 지분상품의 공정가치에 기초하여 거래를 측정하는 경우에는 지분상품의 부여조건을 고려하여 측정기준일 현재 공정가치를 측정한다.
  \item 기업이 거래상대방에게 주식기준보상거래를 현금이나 지분상품발행으로 결제받을 수 있는 선택권을 부여한 경우에는, 부채요소와 자본요소가 포함된 복합금융상품을 부여한 것이다.
\end{choices}
\end{question}
\answer{2}
